\section{Què és la pneumàtica?}

Quan puges a un autobús, la porta s'obre amb un «xiuuu». Aquest soroll és aire que s'escapa d'un circuit pneumàtic. La paraula \emph{pneumàtica} ve del grec \emph{pneûma}, que vol dir \textbf{aire}.

\begin{idea}
La \textbf{pneumàtica} és la tecnologia que utilitza \textbf{aire comprimit} (aire a una pressió més gran que l'atmosfèrica) per transmetre energia i produir \textbf{moviment i força}.

Un \textbf{circuit pneumàtic} és un conjunt d'elements connectats amb tubs que \textbf{produeixen}, \textbf{controlen} i \textbf{aprofiten} l'aire comprimit per fer un treball.
\end{idea}

\subsection*{Les parts d'un circuit pneumàtic}

\begin{center}
\begin{tikzpicture}[
  caixa/.style={draw=#1,fill=#1!10,rounded corners=3pt,minimum width=3.1cm,minimum height=1.5cm,align=center,font=\small,line width=0.8pt},
  fletxa/.style={-{Stealth[length=2.5mm]},line width=1pt,aire}]
 \node[caixa=blau] (a) at (0,0) {\textbf{1. Producció}\\[-1pt]\footnotesize compressor\\[-3pt]\footnotesize i dipòsit};
 \node[caixa=verd] (b) at (3.8,0) {\textbf{2. Condicionament}\\[-1pt]\footnotesize filtre, regulador\\[-3pt]\footnotesize i lubricador};
 \node[caixa=taronja] (c) at (7.6,0) {\textbf{3. Control}\\[-1pt]\footnotesize vàlvules\\[-3pt]\footnotesize (distribuïdores...)};
 \node[caixa=vermell] (d) at (11.4,0) {\textbf{4. Treball}\\[-1pt]\footnotesize actuadors\\[-3pt]\footnotesize (cilindres, motors)};
 \draw[fletxa] (a)--(b); \draw[fletxa] (b)--(c); \draw[fletxa] (c)--(d);
\end{tikzpicture}
\end{center}
\vspace{-0.3em}
{\small L'aire comprimit es \textbf{produeix}, es \textbf{prepara}, es \textbf{dirigeix} cap on cal i finalment es \textbf{transforma en moviment}.}

\begin{sabiesque}
La pneumàtica és a tot arreu: portes d'autobusos i trens, frens de camions, martells pneumàtics, tornavisos i claus d'aire dels tallers, màquines d'embotellar i d'embalar, robots industrials, cadires de dentista, pistoles de pintar\ldots
\end{sabiesque}

\subsection*{Avantatges i inconvenients}

\begin{center}
\renewcommand{\arraystretch}{1.45}
\begin{tabular}{@{}p{7.7cm}p{7.7cm}@{}}
\toprule
\textcolor{verd}{\textbf{Avantatges}} & \textcolor{vermell}{\textbf{Inconvenients}}\\
\midrule
L'aire és \textbf{gratuït i hi ha a tot arreu}. & L'aire és \textbf{compressible}: els moviments no són gaire precisos a velocitats baixes.\\
És \textbf{net}: no embruta (l'aire s'escapa a l'ambient).& Les \textbf{forces} que s'aconsegueixen són limitades.\\
És \textbf{segur}: no fa espurnes ni hi ha risc d'incendi.& Produir aire comprimit \textbf{costa energia}.\\
Moviments \textbf{ràpids} i components senzills i barats.& L'escapament de l'aire fa \textbf{soroll}.\\
\bottomrule
\end{tabular}
\end{center}

\Needspace{14\baselineskip}
\subsection*{Pneumàtica o hidràulica?}

Si en comptes d'aire fem servir un \textbf{líquid} (oli), parlem d'\textbf{hidràulica}. Les dues tecnologies s'assemblen molt, però no fan el mateix.

\begin{center}
\renewcommand{\arraystretch}{1.45}
\begin{tabular}{@{}lp{5.6cm}p{5.6cm}@{}}
\toprule
 & \textbf{Pneumàtica} & \textbf{Hidràulica}\\
\midrule
\textbf{Fluid} & Aire (un gas) & Oli (un líquid)\\
\textbf{Compressibilitat} & Es comprimeix molt & Gairebé no es comprimeix\\
\textbf{Pressió de treball} & Baixa: uns \SIrange{6}{10}{\bar} & Alta: fins a centenars de bar\\
\textbf{Força} & Petita o mitjana & Molt gran (excavadores, premses)\\
\textbf{Velocitat} & Molt alta & Més lenta\\
\textbf{Precisió} & Menor & Major\\
\textbf{Neteja} & Molt neta & Pot perdre oli\\
\bottomrule
\end{tabular}
\end{center}

\subsection*{Exercicis}

\begin{exercici}{1}
Indica si cada màquina és més probable que sigui \textbf{pneumàtica} (P) o \textbf{hidràulica} (H):\\[0.3em]
a) Braç d'una excavadora \quad b) Porta d'un autobús \quad c) Gat per aixecar cotxes \quad d) Martell pneumàtic \quad e) Frens d'un camió \quad f) Premsa que aixafa cotxes vells
\resposta{1.2cm}{a) H; b) P; c) H; d) P; e) P; f) H.}
\end{exercici}

\begin{exercici}{1}
Ordena els elements d'un circuit pneumàtic segons el camí que fa l'aire, numerant-los de l'1 al 4:\\
\underline{\hspace{0.6cm}} Actuador (cilindre) \quad \underline{\hspace{0.6cm}} Compressor \quad \underline{\hspace{0.6cm}} Filtre i regulador \quad \underline{\hspace{0.6cm}} Vàlvula
\resposta{1.2cm}{Compressor (1), filtre i regulador (2), vàlvula (3), actuador (4).}
\end{exercici}

\begin{exercici}{2}
Explica amb les teves paraules per què una excavadora fa servir hidràulica i una porta d'autobús fa servir pneumàtica.
\resposta{3cm}{L'excavadora necessita forces enormes: l'oli no es comprimeix i treballa a pressions molt altes. La porta només necessita una força petita i ràpida, i l'aire és net, barat i segur.}
\end{exercici}

\begin{exercici}{2}
Pensa en una aplicació de la pneumàtica que hagis vist (a casa, al carrer, a l'institut\ldots). Descriu quin moviment fa i quin avantatge té fer-hi servir aire comprimit.
\resposta{3cm}{Resposta oberta. Exemple: les portes de l'autobús; fan un moviment ràpid i net, i són segures perquè no hi ha electricitat a la porta.}
\end{exercici}

\begin{comprova}
Un circuit pneumàtic transforma l'aire comprimit en\ldots\ \ $\square$ llum \quad $\square$ moviment i força \quad $\square$ calor
\resposta{0pt}{En moviment i força.}
\end{comprova}
